\section{Revelation to the Wise}

\begin{quotex}
Revelation presupposes emptiness — space put at its disposal — in order to manifest itself. This is why it is necessary to renounce. \flright{\textsc{Valentin Tomberg}}

The concept of first-born takes on a cosmic dimension. Christ, the incarnate Son, is God's first thought, preceding all creation, which is ordered toward him and proceeds from him. He is both the beginning and the goal of the new creation. \flright{\textsc{Joseph Ratzinger}}

\end{quotex}
\paragraph{Historical Background}
Balaam was a pagan soothsayer, quoted in the Bible. He would have used auguries to come to his prophecies. This is what he prophesied about Israel:

\begin{quotex}
I see him, but not now; I behold him, but not nigh: a star shall come forth out of Jacob and a scepter shall rise out of Israel (Num 24:17) 

\end{quotex}
This prophecy about a star and a king was known to other nations outside Israel. Some of them may have pondered about the new ``king of the Jews" and how to recognize him. The point here is that boundaries were not as strict as we may assume today. A pagan, since he was quoted in the Bible, had a divine revelation. Although Fr. Ratzinger does not mention this, there are parallels between Zoroaster and Moses.\footnote{\url{https://www.jewishencyclopedia.com/articles/15283-zoroastrianism}}

\paragraph{The Magi}
Feast of the Epiphany, January 6. In ``King of the World", \textbf{Rene Guenon} writes:

\begin{quotex}
Now it can be revealed that these mysterious personages represent in truth nothing other than the three leaders of Agarttha. The Mahanga offers gold to Christ and hails him as King; the Mahatma, offering incense, greets him as Priest; and the Brahatma, greeting him as Prophet, proffers myrrh. 

\end{quotex}
The Magi were from the ``land of sunrise", i.e., the East. They were members of the Persian priestly caste, presumably the Zoroastrian religion. The Magi were also the teachers of the Greek philosophers. Aristotle mentioned the Magi, and he believed the Hebrews originated in India where they were called ``Kalani".

On the one hand, the Magi were seekers of truth. On the other hand, they could also be the possessors of supernatural knowledge who used their powers to deceive and seduce (See Acts 13:10). This reflects the difference between sacred magic and personal magic. (See Letter V, \textit{The Pope}). The Magi of Matthew's Gospel are of the former type and use ``religious and philosophical wisdom as an incentive to set off in the right direction, it is the wisdom that ultimately leads people to Christ."

\paragraph{Wisdom}
Besides Baalam's prophecy, according to Tacitus and Suetonius there was speculation at that time that the ruler of the world would emerge from Judah. Fr. Ratzinger speculates about the Magi:

\begin{quotex}
They were people of inner unrest, people of hope, people on the lookout for the true star of salvation. The men of whom Matthew speaks were not just astronomers. They were ``wise". They represent the inner dynamic of religion toward self-transcendence, which involves a search for truth, a search for the true God and hence ``philosophy" in the original sense of the word. Wisdom, then, serves to purify the message of science: the rationality of that message does not remain at the level of intellectual knowledge, but seeks understanding in its fullness and so raises reason to its loftiest possibilities. 

\end{quotex}
Pope Leo XIII, in the Encyclical \emph{On the Nature of Human Liberty} (\emph{Libertas Praestantissimum}), describes the qualities of the wise man:

\begin{quotex}
Even the heathen philosophers clearly recognized this truth, especially they who held that the wise man alone is free; and by the term ``wise man" was meant, as is well known, the man trained to live in accordance with his nature, that is, in justice and virtue. 

\end{quotex}
We have seen revelation in various guises. There was the revelation to Mary who was sinless, to Joseph who was just, and to the Shepherds, who were spiritually poor and watchful. Finally, we have the revelation to the intelligent, beautifully expressed. Fr. Ratzinger shows us the proper use of intelligence. This is how he described the qualities of God:

\begin{quotex}
We may not ascribe to God anything nonsensical or irrational, or anything that contradicts His creation. 

\end{quotex}
So to become more God-like, we need to strive to be sensible, rational, and consistent. \emph{Pace} certain irrationalists and fideists, the true religion is not irrational, although it transcends reason. Rene Guenon asserts nothing different.

Fr. Ratzinger makes some important affirmations. The Magi represent certain types, from which he concludes:

\begin{itemize}
\item They represent the religions moving toward Christ. 
\item They represent the self-transcendence of science toward him. 
\item They are the successors of Abraham who set off on a journey in response to God's call. 
\item They are the successors of Socrates and his habit of questioning above and beyond conventional religion toward the higher truth. 
\end{itemize}
In summary, Christ is not just the fulfillment of the religion of the Hebrews, but also that of the pagans. He is known not just to the priests and prophets, but also to the philosophers and scientists.

\begin{quotex}
The key point is this: the wise men from the east are a new beginning. They represent the journeying of humanity toward Christ. They initiate a procession that continues throughout history. Not only do they represent the people who have found the way to Christ: they represent the inner aspiration of the human spirit, the dynamism of religions and human reason toward him. 

\end{quotex}
\paragraph{The Star}
It is seldom helpful to speculate about likely, or more often not-very-likely, stories to try to harmonize theological accounts with alleged historical and scientific events. So rather than postulating some sort of supernova, there is another way to look at it. In any event, the reliance on a supernatural star destroys the whole inner meaning of the story, which depends on the inner attitude of the Magi, not some external compulsion.

In the year 7 BC (the most probable year for Jesus' birth), there was a conjunction of Jupiter and Saturn in the constellation Pisces (the Fish). Jupiter was the highest deity and would have been brightest in the evening alongside Saturn. Saturn represented the Jewish people, since Saturn was identified with Yahweh (they shared the same day of the week: Saturday). The Babylonian astronomers could conclude this: the birth in the land of the Jews of a ruler who would bring salvation. Of course, this was not compelling in itself, but depended on the Magi's inner state of expectation. Fr. Ratzinger writes:

\begin{quotex}
This implies that the cosmos speaks of Christ. The language of creation provides a great many pointers. It gives man an intuition of the Creator. 

\end{quotex}
Here we can recall Guenon who wrote that Creation is the reflection of the Logos, the first Thought of God. This is the message to the wise. For others:

\begin{quotex}
The knowledge that emerges from creation, and acquires concrete form in the religions, can also become disoriented, so that it no longer prompts man to transcend himself, but induces him to lock himself into systems with which he believe he can, in some way, oppose the hidden powers of the world. 

\end{quotex}
A further point is that astrology, at least as understood, came to an end with the Magi. The ancients regarded the heavenly bodies as divine powers determining men's fate. Now we understand that Christ conquered all the powers and forces in the heavens and reigns over the entire universe. It is not the star that determines the child's destiny, it is the child that directs the star.



\flrightit{Posted on 2015-01-05 by Cologero }

\begin{center}* * *\end{center}

\begin{footnotesize}\begin{sffamily}



\texttt{Dennis on 2017-01-04 at 16:39 said: }

``7 BC (the most probable year for Jesus' birth)". Um no. What is your evidence for this claim? I've never heard a claim for his birth year earlier than 4 BC. I also recommend you read ``Chronicle of the Living Christ" by Robert Powell (translator of Tomberg), which presents a detailed and convincing argument for putting the birth of Christ at December 6, 2 BC.


\hfill

\texttt{Cologero on 2017-01-04 at 22:01 said: }

Thanks for clearing that up for us, Dennis. I am familiar with Powell's argument but disagree with such a purely quantitative approach.

Ultimately, it is not the year that matters, but rather the idea that the birth of Christ is tied to the Winter Solstice, after which ``darkness decreases and light increases". See \textit{The Mystery of Christmas} by Dom Prosper Gueranger for a more complete explanation. (H/T Carman)

\url{http://traditioninaction.org/religious/g009_Mystery.htm}

\end{sffamily}\end{footnotesize}
